\documentclass[aspectratio=169]{beamer}
\usetheme{Madrid}
\usefonttheme{structurebold}
\usefonttheme[onlymath]{serif}

\definecolor{BKTemplate}{rgb}{0.11765, 0.30588, 0.47451}
\usecolortheme[named=BKTemplate]{structure}
\setbeamercolor{title}{fg=BKTemplate,bg=}
\setbeamertemplate{caption}[numbered]

% Engine-adaptive font configuration
\usepackage{iftex}
\ifLuaTeX
  \usepackage{fontspec}
  \setmainfont{Noto Sans}
  \setsansfont{Noto Sans}
  \setmonofont{Noto Sans Mono}
\else\ifXeTeX
  \usepackage{fontspec}
  \setmainfont{Noto Sans}
  \setsansfont{Noto Sans}
  \setmonofont{Noto Sans Mono}
\else
  \usepackage[utf8]{inputenc}
  \usepackage[T5,T1]{fontenc}
\fi\fi

\usepackage[english,vietnamese]{babel}
\usepackage{url}
\usepackage{booktabs}
\usepackage{array}
\usepackage{amsfonts}
\usepackage{amsmath}
\usepackage{mathtools}
\usepackage{xcolor}
\usepackage{colortbl}
\usepackage{graphicx}
\usepackage{tikz}
\usetikzlibrary{shapes,arrows,positioning,calc,decorations.pathreplacing}
\usepackage{listings}

\newcolumntype{L}[1]{>{\raggedright\arraybackslash}p{#1}}

% Code listing configuration
\lstset{
  basicstyle=\ttfamily\fontsize{6.8}{8.2}\selectfont,
  keywordstyle=\color{blue}\bfseries,
  commentstyle=\color{gray!70!black},
  stringstyle=\color{red!70!black},
  breaklines=true,
  frame=single,
  backgroundcolor=\color{BKTemplate!5},
  rulecolor=\color{BKTemplate!40},
  tabsize=2,
  showstringspaces=false,
  columns=fullflexible,
  keepspaces=true
}
\lstdefinestyle{dense}{basicstyle=\ttfamily\fontsize{6.2}{7.2}\selectfont}

% HCMUT backgrounds (skipped silently if the image files are missing)
\newcommand{\setHCMUTFrontBackground}{%
  \setbeamertemplate{background canvas}{%
    \IfFileExists{images/FrontBackground.pdf}{%
      \includegraphics[width=\paperwidth,height=\paperheight]{images/FrontBackground.pdf}%
    }{}%
  }%
}
\newcommand{\setHCMUTInnerBackground}{%
  \setbeamertemplate{background canvas}{%
    \IfFileExists{images/Background.pdf}{%
      \includegraphics[width=\paperwidth,height=\paperheight]{images/Background.pdf}%
    }{}%
  }%
}

% Common TikZ styles
\tikzset{
  arr/.style={->, >=stealth, line width=1.2pt},
  lbl/.style={font=\tiny, align=center, inner sep=1pt},
  hw/.style={rectangle, draw=BKTemplate, fill=blue!10, rounded corners=2pt, align=center, line width=1pt, font=\scriptsize},
  hwred/.style={hw, draw=red!80!black, fill=red!10},
  hwgreen/.style={hw, draw=green!50!black, fill=green!10},
  hworange/.style={hw, draw=orange!80!black, fill=orange!10}
}


%------------------------------------------------------------
% Shared slide helpers
%------------------------------------------------------------
\tikzset{
  key/.style={rectangle, rounded corners=3pt, draw=BKTemplate, fill=BKTemplate!8, line width=0.8pt, align=center, font=\small, inner sep=5pt},
  big/.style={font=\Huge\bfseries, text=red!75!black}
}

\newcommand{\cue}[2]{%
  \par\noindent\begin{tikzpicture}[baseline=-0.65ex]\node[circle, fill=BKTemplate, text=white, font=\scriptsize\bfseries, inner sep=0pt, minimum size=0.42cm]{#1};\end{tikzpicture}\hspace{0.18cm}{\small #2}\par\vspace{0.16cm}}


% Chapter divider
\newcommand{\ChapterDivider}[3]{%
  \begin{frame}[plain]
    \centering\vspace{1.4cm}
    {\Large\textcolor{gray}{#1}}\\[0.25cm]
    {\Huge\textbf{\textcolor{BKTemplate}{#2}}}\\[0.5cm]
    {\normalsize #3}
  \end{frame}}

%------------------------------------------------------------
% Title: edit the course, topic, subtitle, author and date below.
% Thay các placeholder trong dấu ngoặc vuông bằng thông tin bài mới.
%------------------------------------------------------------
\title[Tên đề tài]
{%
  \texorpdfstring{%
    \vspace*{0.8cm}\\
    \textbf{[TÊN MÔN HỌC] ([MÃ MÔN HỌC])}\\[0.22cm]
    \Large\textbf{[TÊN ĐỀ TÀI BÁO CÁO]}%
  }{[TÊN ĐỀ TÀI BÁO CÁO]}%
}
\subtitle{[Mô tả ngắn về nội dung báo cáo]}
\author[Tác giả]
{%
  \texorpdfstring{%
    \textbf{[Họ và tên tác giả / nhóm]}\\
    \textbf{Giảng viên: [Họ và tên]}%
  }{[Họ và tên tác giả / nhóm], Giảng viên: [Họ và tên]}%
}
\date{[Ngày trình bày]}

\begin{document}

{
\setHCMUTFrontBackground
\begin{frame}[plain]
  \titlepage
\end{frame}
}
\setHCMUTInnerBackground

\begin{frame}{Mục lục và thành viên}
  \begin{center}
    \begin{tikzpicture}[r/.style={rectangle, rounded corners=4pt, minimum width=4.2cm, minimum height=2.3cm, align=center, line width=1pt, text width=3.9cm}]
      \node[r, draw=BKTemplate, fill=BKTemplate!8] (a) at (0,0) {{\small\bfseries\color{BKTemplate} Chương 1 -- Nội dung}\\[0.1cm]\footnotesize [Tóm tắt chương 1]};
      \node[r, draw=BKTemplate, fill=BKTemplate!8] (b) at (4.7,0) {{\small\bfseries\color{BKTemplate} Chương 2 -- Công thức}\\[0.1cm]\footnotesize [Tóm tắt chương 2]};
      \node[r, draw=red!80!black, fill=red!8] (c) at (9.4,0) {{\small\bfseries\color{red!75!black} Chương 3 -- Bảng}\\[0.1cm]\footnotesize [Tóm tắt chương 3]};
      \draw[arr] (a) -- (b); \draw[arr] (b) -- (c);
    \end{tikzpicture}
  \end{center}
  \vspace{0.2cm}
  \begin{table}
    \centering
    \small
    \renewcommand{\arraystretch}{1.25}
    \begin{tabular}{l l l l}
      \toprule
      \textbf{Chương} & \textbf{MSSV} & \textbf{Họ và tên} \\
      \midrule
      Chương 1 & [MSSV] & [Họ và tên] \\
      Chương 2 & [MSSV] & [Họ và tên] \\
      Chương 3 & [MSSV] & [Họ và tên] \\
      \bottomrule
    \end{tabular}
  \end{table}
\end{frame}

% Mỗi file chứa tên chương, trang chia chương và slide mẫu.
% Thêm chương tại đây theo thứ tự trình bày.
% Thay tên chương, mô tả và nội dung trong dấu ngoặc vuông.
\section{Chương 1: Nội dung}
\ChapterDivider{Chương 1}{Nội dung}{[Mô tả chương 1]}

\begin{frame}{[Tiêu đề slide nội dung]}
  \begin{block}{[Tiêu đề khối nội dung]}
    [Đoạn giới thiệu vấn đề hoặc ý chính của slide.]
  \end{block}
  \vspace{0.2cm}
  \cue{1}{[Ý chính thứ nhất.]}
  \cue{2}{[Ý chính thứ hai.]}
  \cue{3}{[Thông tin bổ sung hoặc nhận xét.]}
\end{frame}

% Công thức dưới đây chỉ minh họa cú pháp; thay bằng công thức của đề tài mới.
\section{Chương 2: Công thức}
\ChapterDivider{Chương 2}{Công thức}{[Mô tả chương 2]}

\begin{frame}{[Tiêu đề slide công thức]}
  \begin{block}{[Tên công thức]}
    \[
      y = a x + b
    \]
  \end{block}
  \begin{itemize}
    \item $x$, $y$: [Giải thích các đại lượng.]
    \item $a$, $b$: [Giải thích các tham số.]
  \end{itemize}
\end{frame}

% Thay các placeholder bằng tiêu chí và dữ liệu của đề tài mới.
\section{Chương 3: Bảng}
\ChapterDivider{Chương 3}{Bảng}{[Mô tả chương 3]}

\begin{frame}{[Tiêu đề slide bảng]}
  \begin{table}
    \centering
    \small
    \renewcommand{\arraystretch}{1.5}
    \begin{tabular}{l c c}
      \toprule
      \textbf{[Cột thứ nhất]} & \textbf{[Cột thứ hai]} & \textbf{[Cột thứ ba]} \\
      \midrule
      {[Tiêu chí A]} & [Giá trị] & [Giá trị] \\
      {[Tiêu chí B]} & [Giá trị] & [Giá trị] \\
      \bottomrule
    \end{tabular}
  \end{table}
  \vspace{0.2cm}
  \centering
  {\small [Nhận xét về bảng hoặc nguồn dữ liệu.]}
\end{frame}


{
\setHCMUTFrontBackground
\begin{frame}[plain]
  \centering
  \vspace{1.5cm}
  {\Huge\textbf{\textcolor{BKTemplate}{CẢM ƠN THẦY VÀ CÁC BẠN}}}\\[0.4cm]
  {\Large\textbf{\textcolor{BKTemplate}{ĐÃ LẮNG NGHE}}}\\[1.2cm]
  \textbf{Hỏi đáp \& Thảo luận (Q\&A)}\\[0.3cm]
  \textit{[Đơn vị / lớp / nhóm thực hiện]}
\end{frame}
}

\end{document}
